\section{Decomposition: the typed five columns}
\label{sec:decomposition}

The introduction separated the scalar BSD product into five kinds of data:
height/regulator, analytic/period, finite-source, local $p$-adic, and
determinant-assembly data.  We now make this separation into a typed object.
The point is not to replace the Bloch--Kato determinant line by five
unrelated objects, but to keep track of the five native sources before they
are assembled into the common determinant-line target.

Fix an elliptic curve $E/\Q$ and put $M_E=h^1(E)(1)$.  The five columns are
as follows.  The height/regulator column records the Mordell--Weil lattice
and the Neron--Tate determinant.  The analytic/period column records the
leading coefficient and period data.  The finite-source column records
Tamagawa factors, torsion, the $p$-primary $\Sha$ data, and Cassels--Tate
pairings.  For each prime $p$, the local $p$-adic column records the
ordinary or signed local source data.  The determinant-assembly column
records the orientation, sign, and unit conventions needed to compare all
four arithmetic columns in one determinant line.

The five-column map and the common Schur-residual test are summarized in
\Cref{tab:five-column-map,fig:five-column-schematic}.  The table fixes the
paper-prose names used in the later column sections; the figure records the
same bookkeeping as a source-to-target diagram.
\begin{table}[tbp]
\centering
\small
\caption{The five typed columns and their adequacy residuals.  Each row
separates the native source data from the determinant-line target before the
scalar BSD product is read out.}
\label{tab:five-column-map}
\begin{tabularx}{\textwidth}{@{}>{\raggedright\arraybackslash}p{0.18\textwidth}
>{\raggedright\arraybackslash}p{0.25\textwidth}
>{\raggedright\arraybackslash}p{0.22\textwidth}
>{\raggedright\arraybackslash}X@{}}
\toprule
Column & Carrier map & Residual & Adequate when \\
\midrule
Height/regulator &
\(\rho_{\mathrm{ht/reg}}^{\R}\) from the Mordell--Weil lattice and
Neron--Tate pairing &
\(\XiC\) for \(d(\log\det)_H(\delta H)=\tr(H^{-1}\delta H)\) &
the full height pairing is admitted, so the identity source kills the Schur
residual. \\
\addlinespace
Analytic/period &
\(\rho_{\mathrm{an/period}}^{\R}\) from \((A_E,\Omega_E^+)\) &
\(\XiC\) for \(\log|A_E|-\log\Omega_E^+\) (\(A_E\neq0\)) &
both logarithmic coordinates are present; either coordinate alone leaves a
one-dimensional residual. \\
\addlinespace
Finite/Tamagawa &
\(\rho_{\mathrm{finite}}^{\R}\) from Tamagawa, torsion, \(\Sha[p^\infty]\),
and Cassels--Tate data &
\(\XiC\) for the typed finite-source readout &
the full finite-source tower, including the Cassels--Tate pairing, is
retained. \\
\addlinespace
\(p\)-adic &
\(\rho_p\) from ordinary or signed local data &
\(\XiC\) for the ordinary/signed local readout &
the appropriate ordinary or signed local source is admitted, not merely an
unsigned public shadow. \\
\addlinespace
Determinant assembly &
\(\rho_{\det}\) for orientation, sign, and unit compatibility &
\(\XiC\) for the assembled determinant comparison &
the determinant-orientation coordinate is included with the four arithmetic
columns. \\
\bottomrule
\end{tabularx}
\end{table}

\begin{figure}[tbp]
\centering
\begin{tikzpicture}[
  font=\small,
  column/.style={draw, rounded corners=2pt, align=center, minimum width=2.6cm,
    minimum height=1.05cm, fill=black!3},
  residual/.style={draw, rounded corners=2pt, align=center, minimum width=7.5cm,
    minimum height=0.95cm, fill=white},
  target/.style={draw, rounded corners=2pt, align=center, minimum width=4.1cm,
    minimum height=1.05cm, fill=black!8},
  arrow/.style={-{Latex[length=2mm]}, thick}
]
\node[column] (ht) at (-5.6,1.7) {height\\regulator};
\node[column] (an) at (-2.8,1.7) {analytic\\period};
\node[column] (fin) at (0,1.7) {finite\\Tamagawa};
\node[column] (pad) at (2.8,1.7) {\(p\)-adic\\local};
\node[column] (det) at (5.6,1.7) {determinant\\assembly};
\node[residual] (xi) at (0,0) {Schur residual:\quad
  \(\XiC=K_{DD}-K_{DL}K_{LL}^{\dagger}K_{LD}\)};
\node[target] (bk) at (0,-1.55) {Bloch--Kato determinant-line target};

\foreach \n/\shift in {ht/-3cm,an/-1.5cm,fin/0cm,pad/1.5cm,det/3cm}
  \draw[arrow] (\n.south) .. controls +(0,-0.45) and +(0,0.45) ..
    ([xshift=\shift]xi.north);

\draw[arrow] (xi.south) -- (bk.north);
\node[align=center, font=\footnotesize] at (0,-2.6)
  {adequacy means the admitted native source leaves no Schur residual};
\end{tikzpicture}
\caption{The five native source columns are tested before they are assembled in
the common determinant-line target.}
\label{fig:five-column-schematic}
\end{figure}


The Bloch--Kato target is the determinant-line package attached to $M_E$~\citep{Deligne1987}:
the motive, its finite determinant line, the zeta element, and the chosen
trivialization.  The integrated BSD carrier is the corresponding typed
source object: it keeps the height/regulator, analytic/period,
finite-source, local $p$-adic, and determinant-assembly data in their
separate slots.  The first structural map of the paper is the
component-assembly map from this typed source to the determinant-line
target.

\begin{definition}[Bloch--Kato component-assembly map]
\label{def:apparatus:bk-component-map}
For $M_E=h^1(E)(1)$, let
\[
  C_{\mathrm{BSD}}(E)
  =
  \bigl(C_{\mathrm{ht}}(E),C_{\mathrm{an}}(E),C_{\mathrm{fin}}(E),
  \{C_p(E)\}_p,C_{\det}(E)\bigr)
\]
be the integrated BSD carrier with its typed height/regulator,
analytic/period, finite-source, local $p$-adic, and
determinant-assembly slots.  The source-side determinant slot
\(C_{\det}(E)\) is distinct from the Bloch--Kato target below.  Let
\[
  C_{\mathrm{BK}}(E)
  =
  \bigl(M_E,\Delta_f(M_E),z_{\mathrm{BK}}(M_E),\tau_{\mathrm{BK}}(M_E)\bigr)
\]
be the Bloch--Kato determinant-line target.  The component-assembly map is
the direct sum
\[
  \rho_{\mathrm{BK}}
  =
  \rho_{\mathrm{ht}}
  \oplus \rho_{\mathrm{an}}
  \oplus \rho_{\mathrm{fin}}
  \oplus \Bigl(\bigoplus_p \rho_p\Bigr)
  \oplus \rho_{\det}
  \;:\;
  C_{\mathrm{BSD}}(E)\longrightarrow C_{\mathrm{BK}}(E).
\]
Here $\rho_{\mathrm{ht}}$ sends the Mordell--Weil height and regulator
data to the global cohomological determinant side; $\rho_{\mathrm{an}}$
sends central leading-term and period data to the analytic
zeta/trivialization side; $\rho_{\mathrm{fin}}$ sends Tamagawa, torsion,
$\Sha[p^\infty]$, and Cassels--Tate data to the finite arithmetic
determinant side; $\rho_p$ sends the local $p$-adic analytic/Selmer data,
control data, and local normalization to the $p$-primary determinant side;
and $\rho_{\det}$ records the orientation, sign, and unit compatibility in
the same determinant line.
\end{definition}

The definition is a decomposition of the source side, not a theorem about
the truth of BSD.  Its role is to make the columns visible enough that the
adequacy residual of each column can be tested.  Once the columns are
visible, one may compare them through a common determinant-line
trivialization.  In additive coordinates, the comparison defect splits into
one contribution from each column.

\begin{definition}[Bridge-defect equation]
\label{def:apparatus:bridge-defect-equation}
After choosing a comparison trivialization of the determinant-line target,
write $\Delta_{\mathrm{BSD}}^{\mathrm{BK}}$ for the total Bloch--Kato defect
of the BSD bridge.  The bridge-defect equation is the additive
decomposition
\[
  \Delta_{\mathrm{BSD}}^{\mathrm{BK}}
  =
  E_{\mathrm{ht/reg}}
  +E_{\mathrm{an/period}}
  +E_{\mathrm{finite}}
  +\sum_p E_p
  +E_{\det}.
\]
The five summands are respectively the height/regulator defect, the
analytic/period defect, the finite-source defect, the local $p$-adic
defects, and the determinant-assembly defect.  Equivalently, in
multiplicative coordinates the same assertion says that the total component
ratio is the product of the corresponding five component ratios.
\end{definition}

\begin{theorem}[Component-killing]
\label{thm:apparatus:component-killing}
Assume that the five component defects are all legally comparable in the
chosen determinant-line trivialization.  If
\[
  E_{\mathrm{ht/reg}}=0,\qquad
  E_{\mathrm{an/period}}=0,\qquad
  E_{\mathrm{finite}}=0,\qquad
  E_p=0\ \text{for every }p,\qquad
  E_{\det}=0,
\]
then $\Delta_{\mathrm{BSD}}^{\mathrm{BK}}=0$.  Equivalently, if
$\Delta_{\mathrm{BSD}}^{\mathrm{BK}}\neq 0$, then at least one component
defect is nonzero or one of the asserted component comparisons is not
legally available.\footnote{Verified in Lean as
\texttt{SixBirdsBSD.Apparatus.Decomposition.componentKilling}; see App~\ref{app:formalization}.}
\end{theorem}

\begin{proof}
The comparison trivialization puts all five summands in the same additive
defect group.  Hence the bridge-defect equation of
Definition~\ref{def:apparatus:bridge-defect-equation} applies:
\[
  \Delta_{\mathrm{BSD}}^{\mathrm{BK}}
  =
  E_{\mathrm{ht/reg}}
  +E_{\mathrm{an/period}}
  +E_{\mathrm{finite}}
  +\sum_p E_p
  +E_{\det}.
\]
Under the stated hypotheses every summand on the right is zero.  The sum of
zero component defects is zero, so
$\Delta_{\mathrm{BSD}}^{\mathrm{BK}}=0$.

For the contraposition, suppose
$\Delta_{\mathrm{BSD}}^{\mathrm{BK}}\neq0$.  If all five component
comparisons were legally available and every component defect vanished, the
first part of the proof would force
$\Delta_{\mathrm{BSD}}^{\mathrm{BK}}=0$, a contradiction.  Therefore either
some declared component defect is nonzero, or the attempted componentwise
comparison is not legal in the determinant-line trivialization under
discussion.
\end{proof}

\begin{theorem}[Finite scalar factor is a public shadow]
\label{thm:apparatus:finite-scalar-public-shadow}
Let
\[
  C_{\mathrm{fin}}
  =
  C_{\mathrm{Tamagawa}}\oplus C_{\tors}\oplus
  C_{\Sha[p^\infty]}\oplus C_{\mathrm{CT}}
\]
be the typed finite-source tower.  The public scalar finite factor
\[
  \phi_{\mathrm{fin}}(C_{\mathrm{fin}})
  =
  \frac{|\Sha|\prod_\ell c_\ell}{|E(\Q)_{\tors}|^2}
\]
is not injective on this typed tower.  Equivalently, at a fixed prime $p$,
the valuation readout
\[
  q_p
  =
  v_p|\Sha|+v_p\Bigl(\prod_\ell c_\ell\Bigr)
  -2v_p|E(\Q)_{\tors}|
\]
does not determine the typed finite-source datum.  Thus a scalar
finite-factor row cannot be promoted to a finite-source carrier row without
additional source data.\footnote{Verified in Lean as
\texttt{SixBirdsBSD.Apparatus.Decomposition.finiteScalarPublicShadow}; see
App~\ref{app:formalization}.}
\end{theorem}

\begin{proof}
Fix a prime $p$ and consider only the three valuation coordinates appearing
in $q_p$: the $p$-primary part of $\Sha$, the Tamagawa product, and torsion.
The typed finite tower remembers which block carries a valuation, while the
scalar readout only remembers the integer combination
\[
  v_p|\Sha|+v_p\Bigl(\prod_\ell c_\ell\Bigr)
  -2v_p|E(\Q)_{\tors}|.
\]
Now compare two typed finite inputs.  In the first, put the valuation in the
$\Sha[p^\infty]$ block:
\[
  \bigl(v_p|\Sha|,\ v_p(\prod_\ell c_\ell),\
  v_p|E(\Q)_{\tors}|\bigr)=(1,0,0).
\]
In the second, put it in the Tamagawa block:
\[
  \bigl(v_p|\Sha|,\ v_p(\prod_\ell c_\ell),\
  v_p|E(\Q)_{\tors}|\bigr)=(0,1,0).
\]
The two typed inputs are distinct, since the first records a
$p$-primary $\Sha$ contribution and the second records a Tamagawa
contribution.  Nevertheless the scalar valuation is the same:
\[
  q_p=1+0-2\cdot0=1
  \qquad\text{and}\qquad
  q_p=0+1-2\cdot0=1.
\]
Thus the scalar map identifies two different typed finite-source rows.
Moreover the Cassels--Tate datum is absent from the displayed expression
for $q_p$ altogether.  Changing the Cassels--Tate pairing while holding the
three scalar valuation coordinates fixed leaves $q_p$ unchanged.  Hence the
scalar finite factor is a noninjective public shadow of the typed finite
tower.

This is an audit no-go in the sense used throughout the paper: it records
that the scalar finite row alone does not contain enough information to
recover the typed finite-source row.  It is not a claim that no stronger
arithmetic theorem, supplied with additional source data, can identify or
transport the missing finite information.
\end{proof}

\begin{nonclaim}[Decomposition scope]
\label{nonclaim:apparatus:decomposition-scope}
This section does not prove BSD, the Bloch--Kato conjecture, or ETNC~\citep{BurnsFlach2001}.  It
defines the five-column component map, proves the component-killing
consequence of the bridge-defect equation, and records the finite scalar
public-shadow no-go for the typed bridge.  The no-go is an audit result
about what the scalar row fails to determine; it is not a mathematical
non-existence claim about possible richer finite-source identities.  The
external transports from the five native columns to the determinant-line
target remain recognition sources carried separately, not derived theorems
of this decomposition section.
\end{nonclaim}
